\documentclass[12pt,a4paper]{article}
\usepackage{amsmath}
\usepackage{amssymb}
\usepackage{booktabs}
\usepackage{float}
\usepackage{geometry}
\usepackage{listings}
\usepackage{xcolor}
\geometry{margin=2.5cm}

\lstset{
    language=Matlab,
    basicstyle=\small\ttfamily,
    numbers=left,
    numberstyle=\tiny\color{gray},
    stepnumber=1,
    frame=single,
    breaklines=true,
    tabsize=4,
    keywordstyle=\color{blue},
    commentstyle=\color{green!50!black},
    stringstyle=\color{red!70!black},
    showstringspaces=false,
}

\begin{document}

% ============================================================
\section{BJT Modelling}
% ============================================================

The NPN bipolar junction transistors $Q_1$ and $Q_3$ are modelled using the Ebers-Moll equations:
\begin{equation}
    I_C = I_S \left[\left(e^{V_{BE}/V_T} - 1\right)\left(1 - \frac{V_{BC}}{V_A}\right)
          - \frac{1}{\alpha_R}\left(e^{V_{BC}/V_T} - 1\right)\right]
    \label{eq:IC_NPN_full}
\end{equation}
\begin{equation}
    I_B = I_S \left[\frac{1}{\beta_F}\left(e^{V_{BE}/V_T} - 1\right)
          + \frac{1}{\beta_R}\left(e^{V_{BC}/V_T} - 1\right)\right]
    \label{eq:IB_NPN_full}
\end{equation}

The PNP transistor $Q_2$ can be modelled using the same two equations, but since the
base-emitter and base-collector voltages have reversed polarities to the NPN case, the
emitter-base $V_{EB}$ and collector-base $V_{CB}$ voltages are used instead:
\begin{equation}
    I_C = I_S \left[\left(e^{V_{EB}/V_T} - 1\right)\left(1 - \frac{V_{CB}}{V_A}\right)
          - \frac{1}{\alpha_R}\left(e^{V_{CB}/V_T} - 1\right)\right]
    \label{eq:IC_PNP_full}
\end{equation}
\begin{equation}
    I_B = I_S \left[\frac{1}{\beta_F}\left(e^{V_{EB}/V_T} - 1\right)
          + \frac{1}{\beta_R}\left(e^{V_{CB}/V_T} - 1\right)\right]
    \label{eq:IB_PNP_full}
\end{equation}

where $V_T = \frac{kT}{q}$ is the thermal voltage at temperature $T$, given as $26\,\text{mV}$
(the standard value, assuming a temperature of $300\,\text{K}$), $I_S$ is the saturation current,
$V_A$ is the Early voltage, $\alpha_R$ and $\beta_R$ are the common-base and common-emitter
reverse current gains and $\alpha_F$ and $\beta_F$ are the common-base and common-emitter
forward current gains.

Assuming the circuit is appropriately biased, the three transistors operate in the forward-active
region. In this region, the base-emitter voltage is very close to the threshold voltage of the
base-emitter junction (around $0.7\,\text{V}$ for silicon-based transistors such as the ones used
in the circuit), and the magnitude of the collector-emitter voltage is around a couple of volts.
Therefore, for the NPN transistors:
\begin{align*}
    V_{BE} &\approx 0.7\,\text{V} \\
    V_{CE} &\approx 1.5\,\text{V} \\
    \therefore\; V_{BC} &= V_{BE} - V_{CE} \approx -0.8\,\text{V}
\end{align*}

For $V_T = 26\,\text{mV}$, the exponential terms in the Ebers-Moll equations
(\ref{eq:IC_NPN_full}) and (\ref{eq:IB_NPN_full}) will be about:
\begin{align*}
    e^{V_{BE}/V_T} &\approx e^{26.923} = 4.93\times10^{11} \gg 1 \\
    e^{V_{BC}/V_T} &\approx e^{-30.769}    = 4.337\times10^{-14} \approx 0
\end{align*}

Therefore, the contribution of the reverse currents to the total base and collector currents can
be neglected, and the Ebers-Moll equations for the NPN transistors can be simplified to:
\begin{equation}
    I_C = I_S \cdot \left(e^{V_{BE}/V_T} - 1\right) \cdot \left(1 - \frac{V_{BC}}{V_A}\right)
    \label{eq:IC_NPN}
\end{equation}
\begin{equation}
    I_B = \frac{I_S}{\beta_F} \cdot \left(e^{V_{BE}/V_T} - 1\right)
    \label{eq:IB_NPN}
\end{equation}

The same logic applies in the case of the PNP transistor, and equations
(\ref{eq:IC_PNP_full}) and (\ref{eq:IB_PNP_full}) can be simplified to:
\begin{equation}
    I_C = I_S \cdot \left(e^{V_{EB}/V_T} - 1\right) \cdot \left(1 - \frac{V_{CB}}{V_A}\right)
    \label{eq:IC_PNP}
\end{equation}
\begin{equation}
    I_B = \frac{I_S}{\beta_F} \cdot \left(e^{V_{EB}/V_T} - 1\right)
    \label{eq:IB_PNP}
\end{equation}

An important observation is that while in the forward-active region, the base current is only
dependent on the base-emitter voltage.

% ------------------------------------------------------------
\subsection{Determining $\beta_F$}
% ------------------------------------------------------------

When the base-collector voltage is small in magnitude ($|V_{BC}| \approx 0$), the Early effect
can be ignored and the relationship between the collector current and the base current is linear:
\[
    I_C = \beta_F \cdot I_B
\]
This occurs when $V_{BE} \approx V_{CE} \approx 0.7\,\text{V}$. From the DC output
characteristics of the 2N3904 NPN BJT and the 2N3906 PNP BJT given in the laboratory handout,
$\beta_F$ can be approximated for each transistor by noting the value of the collector current at
$V_{CE} \approx 0.7\,\text{V}$ for the characteristic corresponding to the quiescent value of the
base current of the transistor, and dividing the obtained collector current value by the value of
the base current, $\beta_F = I_C / I_B$.

Table~\ref{tab:beta} presents the quiescent values of the base current $I_B$, the base current
of the chosen characteristic, the value of the collector current at $V_{CE} = 750\,\text{mV}$
and the computed common-emitter forward current gain $\beta_F$ for each of the three transistors.
For each transistor, the DC output characteristic corresponding to the base current closest to the
quiescent value was chosen.

\begin{table}[H]
    \centering
    \begin{tabular}{lcccc}
        \toprule
        BJT & Quiescent $I_B$ (mA) & Chosen $I_B$ (mA) & $I_C$ (mA) & $\beta_F$ \\
        \midrule
        $Q_1$ & 2.7   & 2.5 & 114& 45.6\\
        $Q_2$ & 3.9   & 4   & 116& 29\\
        $Q_3$ & 0.035 & 0.5 & 44& 88\\
        \bottomrule
    \end{tabular}
    \caption{Results used for the calculation of $\beta_F$ for transistors $Q_1$, $Q_2$ and $Q_3$.}
    \label{tab:beta}
\end{table}

% ------------------------------------------------------------
\subsection{Determining $V_A$}
% ------------------------------------------------------------

In the forward-active region, the plot of $I_C$ versus $V_{CE}$ is a straight line when $I_B$
is fixed. For fixed values of the base current $I_B$ and the base-emitter voltage $V_{BE}$,
equation (\ref{eq:IC_NPN}) becomes:
\begin{align}
    I_C &= \beta_F \cdot I_B \cdot \left(1 - \frac{V_{BC}}{V_A}\right)
         = \beta_F \cdot I_B \cdot \left(1 - \frac{V_{BE}}{V_A} + \frac{V_{CE}}{V_A}\right)
         \notag \\[4pt]
    I_C &= \beta_F \cdot I_B \cdot \left(1 - \frac{V_{BE}}{V_A}\right)
         + \beta_F \cdot I_B \cdot \frac{V_{CE}}{V_A}
         = h + m \cdot V_{CE}
    \label{eq:VA_derivation}
\end{align}
where $h = \mathrm{const.} = \beta_F \cdot I_B \cdot \left(1 - \dfrac{V_{BE}}{V_A}\right)$ and
$m = \mathrm{slope} = \dfrac{\beta_F \cdot I_B}{V_A}$.
Therefore, the Early voltage is:
\begin{equation}
    V_A = \frac{\beta_F \cdot I_B}{m}
    \label{eq:VA}
\end{equation}

The slope can be determined from the DC output characteristics. Table~\ref{tab:VA} shows the base
current of the chosen characteristic plot, the measured slope and the calculated Early Voltage for
each transistor.

\begin{table}[H]
    \centering
    \begin{tabular}{lccccc}
        \toprule
        Transistor & $I_B$ (mA) & $\Delta y$ (mA) & $\Delta x$ (mV) & Slope $m$ (mA/mV) & $V_A$ (V) \\
        \midrule
        $Q_1$ & 2.5 & 6& 4466.9& $1.343\times10^{-3}$& 84.885\\
        $Q_2$ & 4   & 8& 4333.6& $1.85\times10^{-3}$& 62.703\\
        $Q_3$ & 0.5 & 2& 4666.8& $4.29\times10^{-4}$& 102.564\\
        \bottomrule
    \end{tabular}
    \caption{Results used for the calculation of $V_A$ for transistors $Q_1$, $Q_2$ and $Q_3$.}
    \label{tab:VA}
\end{table}

% ------------------------------------------------------------
\subsection{Determining $I_S$}
% ------------------------------------------------------------

The base current is given by equation (\ref{eq:IB_NPN}):
\[
    I_B = \frac{I_S}{\beta_F} \cdot \left(e^{V_{BE}/V_T} - 1\right)
\]
For any value of the base current, the base-emitter voltage can be determined from the DC input
characteristic of the transistor. Therefore, the only unknown is the saturation current $I_S$,
which can be found by rearranging the equation:
\begin{equation}
    I_S = \frac{I_B \cdot \beta_F}{e^{V_{BE}/V_T} - 1}
    \label{eq:IS}
\end{equation}

Table~\ref{tab:IS} shows the quiescent value of the base current and the corresponding
base-emitter voltage for each transistor, together with the calculated value of $I_S$.

\begin{table}[H]
    \centering
    \begin{tabular}{lccc}
        \toprule
        Transistor & $I_B$ (mA) & $V_{BE}$ (mV) & $I_S$ (A) \\
        \midrule
        $Q_1$ & 2.7   & 779.6& $1.17\times10^{-14}$\\
        $Q_2$ & 3.9   & 788& $7.78\times10^{-15}$\\
        $Q_3$ & 0.035 & 700& $6.25\times10^{-15}$\\
        \bottomrule
    \end{tabular}
    \caption{Values used for the calculation of $I_S$ for transistors $Q_1$, $Q_2$ and $Q_3$.}
    \label{tab:IS}
\end{table}

Table~\ref{tab:params} shows the determined parameters for the transistors used in the circuit.

\begin{table}[H]
    \centering
    \begin{tabular}{lccc}
        \toprule
        Transistor & $\beta_F$ & $V_A$ (V) & $I_S$ (A) \\
        \midrule
        $Q_1$ & 45.6& 84.885& $1.17\times10^{-14}$\\
        $Q_2$ & 29& 62.703& $7.78\times10^{-15}$\\
        $Q_3$ & 88& 102.564& $6.25\times10^{-15}$\\
        \bottomrule
    \end{tabular}
    \caption{Parameters of the transistors used in the circuit.}
    \label{tab:params}
\end{table}

% ============================================================
\section{Diode Modelling}
% ============================================================

Each diode is modelled as a diode described by the Shockley Diode equation in series with a
resistor, as presented in Figure~\ref{fig:diode_model}.
\begin{equation}
    I_D = I_{SD} \cdot \left(\exp\!\left(\frac{V_D}{\eta V_T}\right) - 1\right)
    \label{eq:shockley}
\end{equation}

\begin{figure}[H]
    \centering
    % Circuit: ideal diode + series resistor R_D
    \begin{tabular}{c}
        $+ \quad V_D \quad -$ \\[4pt]
        $\longrightarrow\!\!\!\!\bigcirc\!\!\!\!\longrightarrow\!\!\!\![\,R_D\,]\!\!\!\!\longrightarrow$
    \end{tabular}
    \caption{Diode modelled as a Shockley diode in series with a bulk resistor $R_D$.}
    \label{fig:diode_model}
\end{figure}

$I_{SD}$ is the diode's reverse saturation current, $V_T$ is the thermal voltage ($26\,\text{mV}$
for a temperature of $300\,\text{K}$), and $\eta$ is the ideality factor (typically between 1
and 2). The total voltage across the series combination is the sum of the voltage across the
diode $V_D$ and the voltage across the resistor $V_R$. Ohm's law gives:
\begin{equation}
    V_R = R_D \cdot I_D
    \label{eq:ohm}
\end{equation}
Rearranging equation~(\ref{eq:shockley}) for $V_D$ and substituting:
\begin{equation}
    V = \eta V_T \ln\!\left(\frac{I_D}{I_{SD}} + 1\right) + R_D \cdot I_D
    \label{eq:diode_model}
\end{equation}

% ------------------------------------------------------------
\subsection{Reverse Saturation Current $I_{SD}$}
% ------------------------------------------------------------

When the voltage across the diode is negative enough (less than $-1\,\text{V}$), the diode is
far into the reverse bias region and the current flowing through it is not much larger than the
reverse saturation current $I_{SD}$. The contribution of the bulk resistor to the total voltage
drop is negligible (on the order of pico-volts). Therefore:
\[
    I_D = I_{SD}\cdot\left(\exp\!\left(\frac{-1}{0.7\eta}\right) - 1\right)
        \approx I_{SD}(0 - 1) = -I_{SD}
\]
The value of the reverse saturation current is therefore taken as the magnitude of the current
at $V = -1\,\text{V}$ from Table~1.1:
\[
    I_{SD} = 15.2\,\text{pA}
\]

% ------------------------------------------------------------
\subsection{Curve Fitting}
% ------------------------------------------------------------

The coefficients $\eta$ and $R_D$ are found so that the mean square error between the computed
voltage drop $V$ (given by equation~(\ref{eq:diode_model})) and the measured voltage drop (given
in Table~1.1 from the laboratory handout) is minimised. The model is linear in the two unknowns:
\[
    V_k = \eta \cdot \underbrace{\left[V_T \ln\!\left(\frac{I_k}{I_{SD}} + 1\right)\right]}_{\ell_k}
        + R_D \cdot I_k
\]
Let $(V_1, I_1),\,(V_2, I_2),\,\ldots,\,(V_N, I_N)$ be the $N = 26$ data points selected from
Table~1.1 (every second entry, from $V = -600\,\text{mV}$ to $V = 960\,\text{mV}$; the entry
at $V = -1\,\text{V}$ was used to determine $I_{SD}$ and is excluded). As the lecture notes
state, Fermat's Theorem gives necessary conditions for minimality of the mean square error:
\[
    \frac{\partial\,\text{MSE}}{\partial\,\eta} = 0, \qquad
    \frac{\partial\,\text{MSE}}{\partial R_D} = 0
\]
Setting $\partial\,\text{MSE}/\partial\,\eta = 0$ yields:
\begin{equation}
    p_0 + p_1 \cdot \eta + p_2 \cdot R_D = 0
    \label{eq:normal1}
\end{equation}
where
\[
    p_0 = \sum_{k=1}^{N} -V_k \cdot V_T \ln\!\left(\frac{I_k}{I_{SD}}+1\right), \quad
    p_1 = \sum_{k=1}^{N} \left[V_T \ln\!\left(\frac{I_k}{I_{SD}}+1\right)\right]^2, \quad
    p_2 = \sum_{k=1}^{N} I_k \cdot V_T \ln\!\left(\frac{I_k}{I_{SD}}+1\right)
\]
Setting $\partial\,\text{MSE}/\partial R_D = 0$ yields:
\begin{equation}
    q_0 + q_1 \cdot \eta + q_2 \cdot R_D = 0
    \label{eq:normal2}
\end{equation}
where
\[
    q_0 = \sum_{k=1}^{N} -V_k \cdot I_k, \qquad
    q_1 = \sum_{k=1}^{N} V_T \ln\!\left(\frac{I_k}{I_{SD}}+1\right) \cdot I_k, \qquad
    q_2 = \sum_{k=1}^{N} I_k^2
\]

The MATLAB function m-file \texttt{Diode\_Fit.m} that computes the coefficient values and
solves the normal equations is presented in Figure~\ref{fig:diode_fit}.

\begin{lstlisting}[caption={MATLAB function m-file \texttt{Diode\_Fit.m} that computes the
         normal equation coefficients and solves for the diode parameters $\eta$ and $R_D$.},
         label=fig:diode_fit]
function [eta, R_D] = Diode_Fit(ISD, Vt)
% [eta, R_D] = Diode_Fit(ISD, Vt)
% Fits the diode model V = eta*Vt*ln(I/ISD + 1) + R_D*I to measured
% data using least squares (normal equations), solved with Gauss_Elim.
%
% Input  ISD = reverse saturation current (A)
%        Vt  = thermal voltage (V), typically 26e-3
% Output eta = ideality factor
%        R_D = bulk series resistance (Ohm)
%
% Data: every second entry from Table 1.1 in MP1.pdf (26 points),
%       excluding the -1V entry used to determine ISD.
%
% Version 1: created 29/03/2026, O.O'Friel

% --- Data from Table 1.1 (every second point, SI units) ---
V_D_data = 1e-3 * [-600; -200; -100;  -40;  -20;    0;   20;   50; ...
                     150;  250;  350;  400;  440;  480;  520;  560; ...
                     600;  640;  680;  720;  760;  800;  840;  880; ...
                     920;  960];

I_D_data = [ -14.79e-12; -14.2e-12;  -12.8e-12;  -8.27e-12;  -5.07e-12; ...
               0;          8.27e-12;   30e-12;    421e-12;     4.16e-9; ...
              40.4e-9;   128e-9;      315e-9;    784e-9;       1.94e-6; ...
               4.8e-6;    12e-6;      29.65e-6;   73.6e-6;   180e-6; ...
             440e-6;      1.024e-3;    2.240e-3;   4.480e-3;   8.053e-3; ...
              12.8e-3];

% --- Pre-compute log term (outside loop, no magic numbers) ---
ln_term   = Vt .* log(I_D_data ./ ISD + 1);

% --- Normal equation coefficients (Fermat conditions on MSE) ---
sum_ln_sq = sum(ln_term .^ 2);
sum_I_ln  = sum(I_D_data .* ln_term);
sum_V_ln  = sum(V_D_data .* ln_term);
sum_I_sq  = sum(I_D_data .^ 2);
sum_VI    = sum(V_D_data .* I_D_data);

% --- Assemble and solve 2x2 normal equations system ---
A_NE  = [sum_ln_sq, sum_I_ln; sum_I_ln, sum_I_sq];
b_NE  = [sum_V_ln; sum_VI];
c_fit = Gauss_Elim(A_NE, b_NE);

eta = c_fit(1);
R_D = c_fit(2);

fprintf('=== Diode Curve Fit Results ===\n');
fprintf('  eta = %.4f\n',    eta);
fprintf('  R_D = %.4f Ohm\n', R_D);
end
\end{lstlisting}

The resulting coefficient values are shown in Table~\ref{tab:diode_coeffs}.

\begin{table}[H]
    \centering
    \begin{tabular}{lc}
        \toprule
        Coefficient & Obtained value \\
        \midrule
        $p_0$ & $-4.48054$ \\
        $p_1$ & $2.57168$ \\
        $p_2$ & $1.52125\times10^{-2}$ \\
        $q_0$ & $-2.68845\times10^{-2}$ \\
        $q_1$ & $1.52125\times10^{-2}$ \\
        $q_2$ & $2.55060\times10^{-4}$ \\
        \bottomrule
    \end{tabular}
    \caption{Obtained values for the coefficients in equations~(\ref{eq:normal1})
             and~(\ref{eq:normal2}). Note $p_2 = q_1$ since both equal $\sum I_k \ell_k$.}
    \label{tab:diode_coeffs}
\end{table}

This gives the following $2\times2$ system of linear equations with unknowns $\eta$ and $R_D$:
\begin{equation}
    \begin{cases}
        2.57168\cdot\eta + 1.52125\times10^{-2}\cdot R_D = 4.48054 \\[4pt]
        1.52125\times10^{-2}\cdot\eta + 2.55060\times10^{-4}\cdot R_D = 2.68845\times10^{-2}
    \end{cases}
    \label{eq:diode_system}
\end{equation}
The system is solved using \texttt{Gauss\_Elim.m}. The \texttt{Gauss\_Elim} algorithm
converges to the following diode parameters:
\[
    \eta = 1.7286 \qquad R_D = 2.3037\,\Omega
\]

The system can also be solved using \texttt{Gauss\_Seidel.m}. A preliminary convergence
investigation is required before iterating. The Gauss-Seidel iteration matrix is:
\[
    T_{GS} = -(D + L)^{-1}U
\]
where $D$, $L$, $U$ are the diagonal, strict lower, and strict upper parts of the coefficient
matrix respectively. The spectral radius $\rho(T_{GS}) = \max|\lambda_i(T_{GS})|$ must satisfy
$\rho < 1$ for convergence. For the system as written in~(\ref{eq:diode_system}), diagonal
dominance does not hold in the second row
($2.55060\times10^{-4} \not> 1.52125\times10^{-2}$), so the system is rescaled by
substituting $R_D' = R_D / 100$:
\begin{equation}
    \begin{bmatrix} 2.57168 & 1.52125 \\ 1.52125 & 2.55060
    \end{bmatrix}
    \begin{bmatrix} \eta \\ R_D' \end{bmatrix}
    =
    \begin{bmatrix} 4.48054 \\ 2.68845\times10^{-2} \end{bmatrix}
    \label{eq:diode_scaled}
\end{equation}
The rescaled matrix is strictly diagonally dominant in both rows
($2.57168 > 1.52125$ and $2.55060 > 1.52125$), which guarantees convergence of Gauss-Seidel.
The ideality factor $\eta$ is known to lie in the interval $[1, 2]$ for real silicon diodes:
$\eta = 1$ corresponds to the ideal Shockley limit, while $\eta = 2$ arises when recombination
currents dominate. The initial estimate $\eta_0 = 1.5$ is taken as the midpoint of this range,
placing the starting point centrally within the feasible region without bias toward either
extreme. The bulk resistance $R_D$ is a small parasitic quantity; for small-signal diodes it is
typically of the order of a few ohms. The initial estimate $R_{D,0} = 3\,\Omega$
(i.e.\ $R_{D,0}' = 0.03$) is therefore a physically motivated prior that places the starting
point in the correct order of magnitude. The tolerance is set to $10^{-4}$, which is sufficient
precision for the diode parameters. The maximum number of iterations is set to 20, which is
the standard value used for iterative solvers of this type and is more than sufficient given
the guaranteed convergence. The relaxation parameter $\omega$ is set to 1.

The \texttt{Gauss\_Seidel} algorithm produces a spectral radius of $\rho = 0.352811 < 1$
and converges within 6 iterations to the following solution:
\[
    \eta = 1.7286 \qquad R_D = 0.023054 \times 100 = 2.3054\,\Omega
\]
in close agreement with the Gaussian elimination result ($R_D = 2.3037\,\Omega$), with the
small difference attributable to the finite convergence tolerance of $10^{-4}$.

The MATLAB function m-file that implements Gaussian elimination with partial pivoting is
presented in Figure~\ref{fig:gauss_elim}.

\begin{lstlisting}[caption={MATLAB function m-file \texttt{Gauss\_Elim.m} that implements
Gaussian elimination with partial pivoting.}, label=fig:gauss_elim]
function sol = Gauss_Elim(A, b)
%GAUSS_ELIM Solves a system of linear equations using Gaussian elimination
% with partial pivoting.
%   Gauss_Elim(A, b) solves the system of linear equations described by
%   A * x = b using Gaussian elimination with partial pivoting, where A
%   represents the coefficient matrix and b is the vector of constant terms.
%
%   Parameter A must be an n x n square matrix.
%   Parameter b must be an n x 1 column vector.
%
%   The function returns the solution vector sol as an n x 1 vector.
%
%   Example usage: Consider the following system of linear equations:
%
%   2 * x1 + 1 * x2 = 5
%   4 * x1 + 3 * x2 = 11
%
%   Gauss_Elim([2, 1; 4, 3], [5; 11]) returns the solution vector [2; 1].

% Created by O.O'Friel, student no. 23338256
% on 24/03/2026 for EEEN30150 - Modelling and Simulation

% Version 1 - Gaussian elimination with partial pivoting and error checking

% Validate input A
sz = size(A);

if numel(sz) ~= 2
    error("Gauss_Elim: Parameter 'A' must be a 2-dimensional matrix\n");
end

if sz(1) ~= sz(2)
    error("Gauss_Elim: Parameter 'A' must be a square matrix\n");
end

n = sz(1);

% Validate input b
if ~isequal(size(b), [n 1])
    error("Gauss_Elim: Parameter 'b' must be an n x 1 column vector\n");
end

% Work on copies to avoid modifying inputs
Ac = A;
bc = b;

% Forward elimination with partial pivoting
for k = 1:n-1

    % Find the row with the largest entry in column k from row k downwards
    [~, piv] = max(abs(Ac(k:n, k)));
    piv = piv + k - 1;   % convert to absolute row index

    % Swap rows k and piv in both Ac and bc if required
    if piv ~= k
        Ac([k, piv], :) = Ac([piv, k], :);
        bc([k, piv])    = bc([piv, k]);
    end

    % A zero pivot coefficient implies a singular matrix
    if Ac(k, k) == 0
        error("Gauss_Elim: Division by zero - singular system\n");
    end

    % Eliminate entries below the pivot in column k
    for j = k+1:n
        mult       = Ac(j, k) / Ac(k, k);
        Ac(j, k:n) = Ac(j, k:n) - mult * Ac(k, k:n);
        bc(j)      = bc(j)      - mult * bc(k);
    end

end

% Back substitution
sol = zeros(n, 1);
for k = 0:n-1
    sol(n-k) = (bc(n-k) - Ac(n-k, :) * sol) / Ac(n-k, n-k);
end

end
\end{lstlisting}

The MATLAB function m-file that implements the Gauss-Seidel iterative algorithm with SOR
and spectral radius convergence check is presented in Figure~\ref{fig:gauss_seidel}.

\begin{lstlisting}[caption={MATLAB function m-file \texttt{Gauss\_Seidel.m} implementing the
         Gauss-Seidel iterative method with SOR relaxation and a preliminary spectral
         radius convergence check.}, label=fig:gauss_seidel]
function [sol, count, err] = Gauss_Seidel(A, b, x0, N, tol, omega)
%GAUSS_SEIDEL Solves a system of linear equations using the Gauss-Seidel
% method with relaxation and a preliminary convergence investigation.
%   Gauss_Seidel(A, b, x0, N, tol, omega) solves the system described by
%   A * x = b using the Gauss-Seidel iterative method, where A is the
%   coefficient matrix and b is the vector of constant terms. The algorithm
%   starts from initial guess x0 and attempts to converge to the desired
%   tolerance within N iterations. Before iterating, the spectral radius of
%   the Gauss-Seidel iteration matrix is computed as a convergence check.
%
%   Parameter A     must be an n x n square matrix.
%   Parameter b     must be an n x 1 column vector.
%   Parameter x0    must be an n x 1 column vector (initial guess).
%   Parameter N     must be a positive integer (maximum number of iterations).
%   Parameter tol   must be a positive real number (convergence tolerance).
%   Parameter omega must be in the interval (0, 2) (relaxation factor).
%
%   The function returns the solution vector sol as an n x 1 vector, the
%   number of iterations performed, and the final residual norm.
%
%   Example usage: Consider the following system of linear equations:
%
%   3 * x1 + 1 * x2 = 4
%   1 * x1 + 3 * x2 = 4
%
%   Gauss_Seidel([3,1;1,3],[4;4],[0;0],20,1e-6,1) returns [1; 1].

% Created by O.O'Friel, student no. 23338256
% on 28/03/2026 for EEEN30150 - Modelling and Simulation

% Version 1 - Gauss-Seidel with SOR relaxation and spectral radius check

% Validate input A
sz = size(A);

if numel(sz) ~= 2
    error("Gauss_Seidel: Parameter 'A' must be a 2-dimensional matrix\n");
end

if sz(1) ~= sz(2)
    error("Gauss_Seidel: Parameter 'A' must be a square matrix\n");
end

n = sz(1);

% Validate remaining inputs
if ~isequal(size(b), [n 1])
    error("Gauss_Seidel: Parameter 'b' must be an n x 1 column vector\n");
end

if ~isnumeric(x0) || ~isvector(x0)
    error("Gauss_Seidel: Parameter 'x0' must be a vector\n");
end

if ~isscalar(N) || ~isreal(N) || N <= 1
    error("Gauss_Seidel: Parameter 'N' must be an integer greater than 1\n");
end

if ~isscalar(tol) || ~isreal(tol) || tol <= 0
    error("Gauss_Seidel: Parameter 'tol' must be a positive real number\n");
end

% If omega is outside (0,2), reset to 1 (no relaxation)
if omega <= 0 || omega >= 2
    omega = 1;
end

N = floor(N);

% Preliminary convergence investigation via spectral radius
% Decompose A into diagonal D_mat, strict lower L_mat, strict upper U_mat
D_mat = diag(diag(A));
L_mat = tril(A, -1);
U_mat = triu(A,  1);

% Gauss-Seidel iteration matrix: T_iter = -(D_mat + L_mat)^{-1} * U_mat
T_iter = -(D_mat + L_mat) \ U_mat;
rho    = max(abs(eig(T_iter)));

fprintf('Gauss_Seidel: spectral radius = %.6f\n', rho);

if rho >= 1
    warning("Gauss_Seidel: spectral radius >= 1, convergence not guaranteed\n");
end

% Gauss-Seidel iteration with SOR relaxation
x_cur = x0;

for count = 1:N+1

    % Terminate if maximum iterations reached
    if count == N+1
        fprintf('Gauss_Seidel: Maximum number of iterations reached\n');
        break;
    end

    x_upd = x_cur;   % initialise updated estimate from current values

    % Loop over each unknown
    for i = 1:n
        s_hi  = A(i, i+1:n) * x_cur(i+1:n);   % old values  (j > i)
        s_lo  = A(i, 1:i-1) * x_upd(1:i-1);   % new values  (j < i)
        x_gs  = (b(i) - s_hi - s_lo) / A(i, i);
        x_upd(i) = (1 - omega) * x_cur(i) + omega * x_gs;
    end

    % Update estimate and check convergence
    x_cur = x_upd;

    if norm(A * x_cur - b) < tol
        fprintf('Gauss_Seidel: converged in %d iterations\n', count);
        break;
    end

end

sol = x_cur;
err = norm(A * sol - b);

end
\end{lstlisting}

\end{document}

